\subsection{Overview of proofs and techniques}

We prove \expref{Theorem}{thm:main} by constructing a PCP (for an \cclass{NP}-complete language)
that makes two queries to a proof: one query reads a symbol over an alphabet of
size $q$ and the other reads a symbol
over an alphabet of size $q^6$. The PCP has completeness $1-\zeta$
and soundness at most ${4}/{q}$.
As is standard in the PCP literature, the PCP is obtained by composing the so-called
Outer PCP and Inner PCP.

\paragraph{Inner PCP}
The Inner PCP is essentially a probabilistic
testing procedure that tests whether a given function $f: \F_q^m \mapsto \F_q$
satisfies a desired property. Three types of tests have generally been used depending
on the desired setting of parameters, \eg, the number of queries, type of the
acceptance predicate,
completeness and soundness, size of the proof, etc.
\begin{itemize}
\item The low degree test~\cite{FGLSS, AS, AroraSudan, RazSafra} that tests whether $f$ is a polynomial of low degree.
\item The linearity test that tests whether $f$ is linear (= Hadamard codeword) \cite{ALMSS, Kh-maxclique}.
\item The dictatorship test that tests whether $f$ is a dictatorship, \ie, a function
of the form $f(x)=x_i$ for some coordinate $1 \leq i \leq m$, \eg,~\cite{BGS, Has96, Has97}.
\end{itemize}
If $f$ satisfies the desired property (\ie, being linear, low degree, or dictatorship),
then the test passes with probability
(close to) $1$ and conversely, if the test passes with \emph{reasonable} probability,
$f$ must have a non-trivial agreement with another function $g$ that has the desired property.
The low degree test has been analyzed algebraically~\cite{FGLSS, AS, AroraSudan}
and also combinatorially~\cite{RazSafra}
whereas the linearity and dictatorship tests are often amenable to
Fourier analysis, \eg,~\cite{Kh-maxclique, Has96, Has97, KKMO}.

In this paper, we desire an explicit and \emph{good} trade-off between the alphabet size of the PCP and the
soundness parameter. At the Inner PCP level, this is achieved by a linearity test that combines
the elements of the linearity test and the low degree test. Specifically, given
a function  $f: \F_q^m \mapsto \F_q$, we wish to test that $f$ is linear. The standard
BLR test~\cite{BLR}
checks whether $f(x+y)=f(x)+f(y)$ for randomly chosen inputs $x$ and $y$. We
however wish to have a $2$-query test and hence we instead do a \emph{point versus subspace}
test. We are given the table of values of the function $f$ and in addition, for
every subspace $W \subseteq \F_q^m$ of dimension $6$, a linear function $\calT(W): W \mapsto \F_q$
that is
supposed to be the restriction of $f$ on $W$ (denoted $f|_W$). The test selects a random $6$-dimensional
subspace $W$ and a random point $w \in W$ and accepts if and only if
$f(w) = \calT(W) (w)$. Note that the query $f(w)$ is over an alphabet of size $q$ and the
query $\calT(W)$ is over an alphabet of size $q^6$.



We achieve the following soundness guarantee: if $f$ passes the \emph{point versus subspace}
test with probability ${3}/{q}$, then for some $j \in \F_q, j \not= 0$,
the function $j\cdot f$ has a Fourier coefficient with magnitude at least ${1}/{q^2}$
(see \expref{Lemma}{lemma:inner}). Interestingly, the test is analyzed by looking at the
probability that $f$ (and its restriction $f|_W$) passes the \emph{Gowers Test}
$f(x) - f(y) - f(z) + f(-x+y+z)=0$ for randomly chosen $x,y,z$ (see \expref{Section}{sec:gowers-test} for the reason the test is named so).
The Gowers Test is considered
for the purposes of the analysis only and is not a part of the actual test. The probability
of passing the Gowers Test is related to the existence of a \emph{large} Fourier coefficient
(see \expref{Lemma}{lemma:Gowers}) for function $j\cdot f$ for some $j \not= 0$.

The analysis proceeds as follows: assume that $f$ passes the \emph{point versus subspace}
test with probability ${1}/{q}+\delta$ where $\delta \geq {2}/{q}$. For the sake of
simplicity, assume that for every subspace $W$,
the test passes with probability ${1}/{q}+\delta$ after selecting $W$. Thus
$f|_W$ has agreement ${1}/{q}+\delta$ with a linear function $\calT(W)$.
We show that
this implies $j \cdot f|_W$ has a \emph{large} Fourier coefficient for some $j \not= 0$
and hence
$f|_W$ passes the Gowers Test with probability
${1}/{q}+\delta^4$. We observe that the probability that $f$ passes the Gowers Test is the average (over $W$)  of the probability that $f|_W$ passes the Gowers Test, up to an additive difference of
$e = {3}/{q^4}$. This implies that $f$ passes the Gowers Test with probability
${1}/{q} + \delta^4 - e \geq {1}/{q} + {2}/{q^4}$. {}From this we conclude that
for some $j \not= 0$, $j \cdot f$ has a \emph{large} Fourier coefficient, with magnitude
at least ${1}/{q^2}$.

Thus the Gowers Test serves as a vehicle to pass
from the local linearity of $f$ to its global linearity. This is also reminiscent of the
\emph{bootstrapping} method that allows the analysis of the low degree test in two (or three)
dimensions to carry over to a higher number of dimensions.

\paragraph{Remark} It is not clear that we necessarily need a \emph{point
versus $\ell$-dimensional subspace} test with $\ell=6$
to achieve a soundness of $O({1}/{q})$. Here is the
limitation of our current analysis.  In the Gowers Test
on $f|_W$, $dim(W)=\ell$, the
inputs $x, y, z$ are linearly dependent with probability
$\Theta({1}/{q^{\ell-2}})$. On the other hand, when the Gowers Test is
applied to the function $f$, the inputs $x,y,z$ are linearly dependent with probability
$\Theta({1}/{q^{m-2}} )$ which is negligible.  Hence we are able to claim that
the probability that $f$ passes the Gowers Test is the average (over $W$)  of the probability that $f|_W$ passes the Gowers Test, but only up to an additive difference of $e = \Theta({1}/{q^{\ell-2}})$. As the above calculation shows, we desire that $\delta
=O({1}/{q})$ and that $\delta^4$ dominates
$e = \Theta({1}/{q^{\ell-2}})$, which forces us to have $\ell \geq 6$.



\paragraph{Outer PCP, Composed PCP and Sub-Code Covering Property}
The linearity testing primitive at the Inner PCP level dictates that we use an Outer PCP
based on a \cclass{NP}-hard problem with linear constraints. A natural choice is the 3LIN problem
over $\F_q$: we are given an instance $(X, \Phi)$ where $X$ is a set of variables taking
values in $\F_q$ and $\Phi$ is a set of linear equations, each equation depending on three
variables from $X$. The goal is to find an assignment $\sigma: X \mapsto \Phi$ that
satisfies a good fraction of the equations. A well-known result of H\aa stad~\cite{Has97}
shows that for any constant $\eta > 0$,
it is \cclass{NP}-hard to distinguish whether an instance $(X,\Phi)$ has an assignment
that satisfies $1-\eta$ fraction of the equations (YES Case) or any assignment satisfies at most
${1}/{q}+\eta$ fraction of the equations (NO Case).

Starting with the hard instance of 3LIN as above, one builds a 2P1R Game as follows:
the first prover is sent a set of $k$ equations at random from $\Phi$. Let $V$ denote
the set of $3k$ variables sent to the first prover. The second prover is sent a set
$U \subseteq V$ that includes independently for $1 \leq i \leq k$, all three variables
in the $i^{th}$ equation with probability $1-\beta$ and exactly one of the three variables
in the $i^{th}$ equation with probability ${\beta}/{3}$ each
($\beta$ will be tiny as explained later).
The provers answer with assignments to $V$ and $U$ respectively and the verifier accepts
if and only if the assignments are consistent on $U$ and moreover all equations on $V$
are satisfied. In the YES Case, the provers have a strategy to make the verifier
accept with probability at least $1-k\eta$ whereas in the NO Case, any prover
strategy makes the verifier accept with probability at most $2^{-\Omega(\beta k)}$ (see
\expref{Section}{sec:outer} for a proof).

The 2P1R Game described is precisely the so-called Outer PCP\@. The composition of the
Outer and Inner PCP amounts to constructing a verifier that behaves as follows: the
PCP verifier expects, for each question $V$ ($U$ resp.) to the first (second resp.)
prover, a Hadamard encoding of assignment to $V$ ($U$ resp.). The Hadamard code is
same as the table of values of a linear function $f_V : \F_q^V \mapsto \F_q$ ($g_U: \F_q^U \mapsto \F_q$ resp.)
defined by an assignment to $V$ ($U$ resp.). Moreover, for every
$V$, a table of linear functions on all $6$-dimensional subspaces of $\F_q^V$ is
expected; these linear functions are supposed to be the restrictions of the global
linear function on $\F_q^V$. The verifier now picks a random question $V$ to the
first prover and performs the \emph{point versus subspace} test on the table
$f_V$ and the corresponding subspaces table.


Note that it appears as if the Hadamard codes on $U$ do not play any role (which would
not make sense). We observe that the Hadamard code on $U$ is actually contained in
the Hadamard code on $V$ (as a sub-code). This is because $\F_q^U \subseteq \F_q^V$ where
we append a vector in $\F_q^U$ with zeroes at positions in $V \setminus U$
and think of it as a vector in  $\F_q^V$. Thus there is no need to have a separate
Hadamard code on $U$ (though it helps in the analysis to think of these as virtual
tables). Moreover if $U \subseteq V \cap V'$ for distinct questions $V$ and $V'$ to the
first prover, we can identify the positions in Hadamard codes of $V$ and $V'$ that
correspond to the same position in the (virtual) Hadamard code of $U$.\footnote{The
reader might be concerned here that the 3LIN instance has imperfect
completeness and whether such identification can be done. This is a subtle and valid point:
a short answer is that
the identification can be done, but as mentioned in the next paragraph, the standard trick of \emph{folding over equations}
is no longer allowed and the \emph{folding} needs to be implemented explicitly.}

Now we look carefully at the soundness analysis of the composed PCP\@. Assume on the
contrary that the verifier
accepts with probability ${4}/{q}$ and for simplicity that for every $V$,
the verifier accepts
 the \emph{point versus subspace} test on the supposed Hadamard code
$f: \F_q^V \mapsto \F_q$ with probability ${4}/{q}$. The analysis of the Inner PCP
guarantees that for some $j \not=0$, $j\cdot f$ has a large Fourier coefficient.
Thus the function $f$ may be list decoded by making a list of all large Fourier
coefficients. Since the sum of squared magnitudes of all Fourier coefficients is $1$, the list size
is bounded. Our test also incorporates \emph{side-conditions} (see \expref{Section}{sec:point-subspace}) and ensures that
 the Fourier coefficients
obtained as list decoding satisfy all the equations on $V$ (this is done
in a more explicit manner than the standard \emph{folding over equations}
trick which is
inapplicable in our setting).

Finally we want to infer consistency between $f$ (\ie, supposed Hadamard code
on $V$) and the (virtual) supposed Hadamard codes $g_U: \F_q^U \mapsto \F_q$
on $U \subseteq V$. But as observed $g_U = f|_{\F_q^U}$.
We would like to conclude that since $j \cdot f$ has a large Fourier coefficient,
so do \emph{many}  of the $j \cdot g|_U$
functions. This turns out to be  possible if the sub-code spaces
$\F_q^U$ over all choices of $U$ (weighted according to the distribution on $U$ for
a fixed $V$) cover the global space $\F_q^V$ almost uniformly. We term this as
the \emph{sub-code covering property}.
When $\beta$ is sufficiently small and $k$ is
sufficiently large, we note that 
\[
|U| \approx \left(1-\frac{2}{3}\beta\right) |V|\,;
\]
thus the size of
a typical sub-code space is $q^{-O(\beta k)} = 2^{-O(\beta \log q \cdot k)}$
relative to the size of the code, there are roughly 
\[
\binom{k}{\beta  k} \cdot 3^{\beta k}
= 2^{\Omega(\beta \log(1/\beta)\cdot k)}
\]
choices of $U$, and it is
not unreasonable to expect that
the covering property holds provided $\log(1/\beta) \gg \log q$.
We formally prove this as \expref{Lemma}{lemma:covering}.

Once we are able to infer the consistency of tables $f = f_V$ and $g_U$,
as usual, the list decoding and then picking a random Fourier coefficient in the
list yields a provers' strategy in the 2P1R Game. This yields a contradiction
provided the soundness $2^{-\Omega(\beta k)}$ of the 2P1R Game
is low enough, which follows if
$\beta k$ is large enough.

\paragraph{The Codeword Test and the Consistency Test} In the PCP
literature, the low degree test, linearity test and the dictatorship test at the
Inner PCP level
are
often referred to as the \emph{codeword test} and then the Inner/Outer PCP composition
amounts to extending the test to the \emph{consistency test}
between two (or more as is necessary in some PCPs)
supposed codewords, \eg, $f_V$ and $g_U$ as above. Often this composition
presents serious technical challenges and one needs to carefully analyze
each PCP construction by itself. Our paper shows how to translate
the codeword test for Hadamard code (\ie, the linearity test), essentially in a
black-box manner, to a consistency test. In fact the tables $g_U$ are
virtual and there is no separate
consistency test. As described above, the codeword test for functions $f_V$
automatically serves as a consistency test between $f_V$ and $g_U$, since the
entries in tables
$f_V$ and  $f_{V'}$ that correspond to
the virtual table $g_U$ with $U \subseteq V \cap V'$
are identified together. We state this observation as an informal theorem.\footnote{We remark that a similar informal theorem also holds for the dictatorship test
modulo the Unique Games Conjecture, with the notion of Fourier coefficients replaced
by influences of co-ordinates.}
